% Part: sets-functions-relations
% Chapter: arithmetization
% Section: rationals

\documentclass[../../../include/open-logic-section]{subfiles}

\begin{document}

\olfileid[fa-IR]{sfr}{arith}{rat}
\olsection{از $\Int$ به $\Rat$}

دیدیم که چگونه می‌توان با بهره‌گیری از قدری نظریهٔ طبیعی مجموعه‌ها،
اعداد صحیح را از اعداد طبیعی ساخت. اکنون خواهیم دید که چگونه می‌توان
اعداد گویا را به روشی بسیار مشابه از اعداد صحیح ساخت. دو مشاهدهٔ
نخستین ما چنین‌اند:
\begin{enumerate}
	\item هر عدد گویا را می‌توان به‌شکل $\nicefrac{i}{j}$ نوشت،
	که در آن $i$ و $j$ هر دو عدد صحیح‌اند، اما $j$ ناصفر است.
	\item اطلاعات نهفته در عبارت $\nicefrac{i}{j}$ را به‌همان اندازه
	می‌توان با زوج مرتب $\tuple{ i, j}$ رمزگذاری کرد.
\end{enumerate}
رویکرد آشکار آن بود که اعداد گویا را \emph{همان} زوج‌های مرتبِ برگرفته
از $\Int \times (\Int \setminus \{0_\Int\})$ در نظر بگیریم. اما مانند
گذشته، این رویکرد اندکی بیش از اندازه ساده‌انگارانه است، زیرا می‌خواهیم
$\nicefrac{3}{2} = \nicefrac{6}{4}$، درحالی‌که $\tuple{ 3, 2}\neq \tuple{ 6,
4}$. به‌طور کلی‌تر، خواستهٔ ما چنین است:
\[
	\nicefrac{a}{b} = \nicefrac{c}{d} \text{ iff } a \times d = b \times c
\]
برای دستیابی به این خواسته، رابطهٔ هم‌ارزیِ زیر را روی $\Int \times
(\Int \setminus \{0_\Int\})$ تعریف می‌کنیم:
\[
	\tuple{ a, b }\Ratequiv \tuple{ c, d} \text{ iff }a \times d = b \times c
\]
باید بررسی کنیم که این یک رابطهٔ هم‌ارزی است. این بررسی بسیار شبیه
حالت $\Intequiv$ است و آن را به‌عنوان تمرین وامی‌گذاریم.
\begin{prob}
نشان دهید $\Ratequiv$ یک رابطهٔ هم‌ارزی است.
\end{prob}
این رابطه به ما اجازه می‌دهد بگوییم:
\begin{defn}
اعداد گویا، رده‌های هم‌ارزیِ زوج‌های اعداد صحیح تحت $\Ratequiv$
هستند (زوج‌هایی که مؤلفهٔ دومشان ناصفر است). یعنی $\Rat =
\equivclass{(\Int \times (\Int\setminus \{0_\Int\}))}{\Ratequiv}$.
\end{defn}

همانند اعداد صحیح، می‌خواهیم برخی اعمال بنیادی را نیز تعریف کنیم.
هرگاه $\equivrep{i,j}{\Ratequiv}$ ردهٔ هم‌ارزیِ تحت $\Ratequiv$ باشد
که $\tuple{i, j}$ را به‌عنوان یک !!a{element} در بر دارد، می‌گوییم:
\begin{align*}
	\equivrep{a, b}{\Ratequiv} + \equivrep{c, d}{\Ratequiv} &= \equivrep{ad + bc,  bd}{\Ratequiv}\\
	\equivrep{a, b}{\Ratequiv} \times \equivrep{c, d}{\Ratequiv} &= \equivrep{a  c, b d}{\Ratequiv}.
\intertext{برای تعریف $r \leq s$ روی این اعداد گویا، از این واقعیت بهره می‌بریم که
$r \le s$ اگر و تنها اگر $s - r$ منفی نباشد؛ یعنی $r - s$ را بتوان به‌شکل
$\nicefrac{i}{j}$ نوشت، که در آن $i$ نامنفی و $j$ مثبت است:}
	\equivrep{a, b}{\Ratequiv} \leq \equivrep{c, d}{\Ratequiv} &\text{ iff }
	\equivrep{c, d}{\Ratequiv} - \equivrep{a, b}{\Ratequiv} =
	\equivrep{i_\Int, j_\Int}{\Ratequiv}
\end{align*}
برای بعضی $i \in \Nat$ و $0 \neq j \in \Nat$.

سپس باید بررسی کنیم که این تعریف‌ها چنان‌که \emph{باید} رفتار
می‌کنند؛ این بررسی را به \olref[check]{sec} وامی‌گذاریم. اما آن‌ها
به‌راستی چنین‌اند!{} در پایان، می‌خواهیم راهی داشته باشیم تا اعداد
صحیح را \emph{به‌منزلهٔ} اعداد گویا در نظر بگیریم؛ پس برای هر
$i \in \Int$ مقرر می‌کنیم که $i_\Rat =
\equivrep{i, 1_\Int}{\Ratequiv}$. در \olref[check]{sec} دوباره بررسی
می‌کنیم که همهٔ این تعریف‌ها درست رفتار می‌کنند.

\begin{prob}
نشان دهید $(i + j)_\Rat = i_\Rat+ j_\Rat$ و $(i \times j)_\Rat =
i_\Rat \times j_\Rat$ و $i \leq j \liff i_\Rat \leq j_\Rat$، برای هر
$i, j \in \Int$.
\end{prob}

\end{document}
